% ============================================================
% Task 4 section — updated with a Data Cleaning subsection.
% Replace everything from "\section{Task 4: Exploring..." to just
% before "\newpage \section{Task 5:}" with the code below.
% (Also delete the stray empty "\section{Task 4:}" above it.)
% Requires figure produced by Project_2_Question_4.Rmd:
%   task4_los_hist.pdf
% ============================================================

\section{Task 4: Exploring the Hospital Length-of-Stay Data}

\subsection{Data Overview}
The second engagement concerns hospital capacity planning. The data set
(\texttt{mimic3d.csv}, derived from the MIMIC-III critical-care database)
contains \textbf{58{,}976 hospital admissions and 28 variables}. Each record is
one admission and combines patient demographics, admission details, measures
of how intensively the hospital interacted with the patient, and our target
variable: \texttt{LOSdays}, the length of stay in days.
Table~\ref{tab:mimic_dictionary} summarizes the main variable groups.

\begin{table}[H]
\centering
% Reduce the overall font size of the table to save space
\small

% Set up the alternating row rules:
% \rowcolors{<start_row>}{<odd_row_color>}{<even_row_color>}
\rowcolors{2}{gray!10}{white}

\begin{threeparttable}
\resizebox{\columnwidth}{!}{%
\begin{tabular}{p{3.2cm} p{2.3cm} p{7cm}}
    \toprule
    % Override the automatic color for the header row
    \rowcolor{gray!30}
    \textbf{Variable(s)} & \textbf{Data Type} & \textbf{Description} \\
    \midrule
    \texttt{LOSdays} & Numeric & Length of the hospital stay in days (the target variable). \\
    \texttt{gender}, \texttt{age}, \texttt{insurance}, \texttt{ethnicity}, \dots & Categorical / Integer & Patient demographics. Age is capped at 89 for de-identification; newborns are recorded as age 0. \\
    \texttt{admit\_type}, \texttt{admit\_location}, \texttt{AdmitDiagnosis} & Categorical & Circumstances of the admission (elective, emergency, urgent, or newborn) and the admitting diagnosis. \\
    \texttt{NumLabs}, \texttt{NumRx}, \texttt{NumNotes}, \dots & Numeric & Intensity of hospital--patient interactions (laboratory tests, prescriptions, clinical notes, transfers, \dots), expressed \emph{per day of stay}. \\
    \texttt{ExpiredHospital} & Binary (0/1) & In-hospital mortality flag (9.9\% of admissions). \\
    \bottomrule
\end{tabular}%
}
\caption{Key variables in the length-of-stay data set}
\label{tab:mimic_dictionary}
\end{threeparttable}
\end{table}

\subsection{Data Cleaning}
As in Task~1, the data were audited before any inference. The admission
identifier \texttt{hadm\_id} contains no duplicates, but three issues
required attention.

\begin{itemize}[leftmargin=*, itemsep=4pt]
  \item \textbf{Missing demographic values.} Genuine missing entries are
  confined to \texttt{marital\_status} (10{,}128 records), \texttt{religion}
  (458) and \texttt{AdmitDiagnosis} (25). Because these variables are
  categorical, no records were discarded: the missing entries were merged
  into the ``unknown'' categories that already exist in the data
  (\texttt{UNKNOWN (DEFAULT)} for marital status, \texttt{NOT SPECIFIED}
  for religion).
  \item \textbf{Hidden missing values.} \texttt{AdmitProcedure} encodes
  missingness as the literal text \texttt{"na"} in 8{,}032 records, which a
  standard missing-value check silently overlooks. These were recoded as
  true missing values, interpreted as ``no procedure recorded at
  admission''.
  \item \textbf{Zero-day stays.} 128 admissions record a length of stay of
  exactly zero days, and 95 of them (74\%) correspond to patients who died
  in hospital: these are stays too short to register rather than usable
  measurements. Since they cannot inform a model of stay duration -- and
  the logarithmic transformation used below is undefined at zero -- they
  were removed.
\end{itemize}

After cleaning, \textbf{58{,}848 admissions remain} (99.8\% of the original
data). All results that follow use the cleaned data.

\subsection{Length of Stay}
Table~\ref{tab:los_summary} summarizes the distribution of \texttt{LOSdays}.

\begin{table}[H]
\centering
\caption{Summary statistics of length of stay (days, cleaned data)}
\label{tab:los_summary}
\begin{threeparttable}
\begin{tabular}{r r r r r r r}
\toprule
\textbf{Mean} & \textbf{Median} & \textbf{Std.\ Dev.} & \textbf{Min} & \textbf{Max} & \textbf{Skewness} & \textbf{P(LOS \(>\) 2)} \\
\midrule
10.14 & 6.46 & 12.46 & 0.04 & 294.63 & 4.55 & 89.8\% \\
\bottomrule
\end{tabular}
\end{threeparttable}
\end{table}

Three features matter for the modelling in Task~5. First, the distribution is
\textbf{heavily right-skewed} (skewness $\approx 4.6$): the median stay is
about 6.5 days, but a long tail reaching 295 days pulls the mean up to
10.1 days (Figure~\ref{fig:los_hist}). Second, \textbf{89.8\% of admissions
already exceed two days}, a useful empirical baseline for the probability
question posed in Task~5. Third, on the logarithmic scale the skewness falls
from $4.6$ to $-0.5$ and the distribution becomes approximately normal,
\textbf{justifying a normal likelihood with normal priors on
\(\log(\texttt{LOSdays})\)} in Task~5.

\begin{figure}[H]
\centering
\includegraphics[width=0.9\linewidth]{task4_los_hist.pdf}
\caption{Distribution of length of stay (axis truncated at 60 days; the tail
extends to 295 days). The solid line marks the median, the dashed line the mean.}
\label{fig:los_hist}
\end{figure}

\subsection{Length of Stay Across Admission Types}
Table~\ref{tab:los_by_type} breaks the stays
down by admission type. Emergency admissions dominate the data (71\% of
records). Urgent admissions stay longest (median 8.1 days). Newborns are a
distinctive subgroup: they show the \emph{shortest median} (4.1 days) but a
\emph{high mean} (11.5 days), because most newborns go home quickly while a
minority of intensive-care cases stay for weeks. Gender, by contrast, makes
essentially no difference (median $\approx$ 6.5 days for both women and men).

\begin{table}[H]
\centering
\caption{Length of stay by admission type (cleaned data)}
\label{tab:los_by_type}
\begin{threeparttable}
\begin{tabular}{l r r r}
\toprule
\textbf{Admission Type} & \textbf{Admissions} & \textbf{Median LOS} & \textbf{Mean LOS} \\
\midrule
Newborn   & 7{,}841  & 4.08 & 11.53 \\
\addlinespace
Elective  & 7{,}702  & 6.29 & 8.87 \\
\addlinespace
Emergency & 41{,}972 & 6.88 & 10.04 \\
\addlinespace
Urgent    & 1{,}333  & 8.08 & 12.30 \\
\bottomrule
\end{tabular}
\end{threeparttable}
\end{table}
